\begin{statement}{A.33}{Proved}{The nonstable symmetric genus-three covers at two and their odd-prime reduction}
For $a=6$ or $-6$, let
\[
 C_a:\quad v^2=w^7+aw^4+w,\qquad
 X_a=C_a\setminus\{P_0,P_\infty\}.
\]
Use the lift of any original packet of Part A at a selected
cusp, and any permitted open subgroup $\Delta_0\subset\Delta$
preserving this cover and that lift. Then the full pro-$2$
fixed subgroup is exactly the selected cusp inertia:
\[
 \pi_1(X_{a,\overline K})^{(2),\,s_b(\Delta_0)}
                       =I_{a,2}.                       \tag{A.33.1}
\]
For every odd prime $\ell$, every leading ordinary
lower-central class of a fixed element belongs, after the
geometric identifications in the proof, to
\[
 \mathbb L_{\mathbf Q_\ell}(t,\Omega)
 \ \subset\
 \mathbb L_{\mathbf Q_\ell}
       (\,\mathbf Q_\ell t\oplus V_1\oplus V_2\oplus V_3\,),
 \qquad \Omega=\sum_{i=1}^3[a_i,b_i].                  \tag{A.33.2}
\]
Here $t$ is the primitive homology class of the chosen
cusp, and $a_i,b_i$ are symplectic generators of the three
elliptic components at two. Their ordinary degrees are
$\deg t=1$ and $\deg\Omega=2$.
If the actual local commutator has nonzero first Johnson
degree, the odd-prime fixed subgroup is also exactly
the chosen inertia. Otherwise (A.33.2) is a reduction
of the remaining problem; it does not assert the
existence of an extra fixed element.
Statement A.38 below completes this remaining case
by upgrading to the weight filtration and applying
the cyclotomic fixed-subgroup criterion.
\end{statement}
\begin{proof}
\paragraph{The unramified double cover settles the prime two.}
The fixed-point-free reflection is
\[
 j(w,v)=(1/w,-v/w^4).
\]
Its quotient map is
\[
 t_D=w+1/w,\qquad
 z_D=(w-1/w)v/w^2,
\]
and gives the actual curve of A.32:
\[
 z_D^2=(t_D^2-4)(t_D^3-3t_D+a).
\]
The identities follow by substitution. A fixed point
of $j$ away from the cusps would have $w=\pm1$ and
$v=0$. The two values of the polynomial are $a+2$
and $a-2$, both nonzero. The two cusps are exchanged,
so there are no remaining fixed points. Thus
$C_a\to D_a$ is an unramified double cover of proper
curves, and $X_a\to Y_a$ is its restriction.
The two cusps map unramifiedly to the unique cusp
of $Y_a$. The chosen tangential base point therefore
maps to the chosen tangential base point there.

Write $H_C\triangleleft H_D$ for the corresponding
geometric profinite groups. Their quotient is $C_2$.
The elementary core argument in A.22, applied with
ambient group $H_D$, proves that $H_C^{(2)}$ embeds
as an open subgroup of index two in $H_D^{(2)}$.
This assertion does not use any CM hypothesis:
an open normal subgroup with $2$-group quotient has
its pro-$2$ topology induced from the ambient group.
The embedding is equivariant after a permitted
finite extension. A.32 consequently gives
\[
 (H_C^{(2)})^{s_b(\Delta_0)}
       =H_C^{(2)}\cap (H_D^{(2)})^{s_b(\Delta_0)}
       =H_C^{(2)}\cap I_{D,2}=I_{a,2}.
\]
The last equality uses ramification index one at
the selected cusp. If an auxiliary extension was
needed, prove this first for its smaller open
subgroup of $\Delta_0$, and then use the inclusions
of fixed subgroups and the fact that $\Delta_0$
fixes the selected inertia. This proves (A.33.1)
for the originally stated $\Delta_0$.
For odd $\ell$, the double cover instead induces
a surjection on maximal pro-$\ell$ quotients;
the preceding embedding argument is unavailable.

\paragraph{The pointed three-tail reduction at two.}
The cyclic cubic quotient is
\[
 f:C_a\longrightarrow E_{a,1},\qquad
          (w,v)\longmapsto(w^3,wv),
 \qquad E_{a,1}:\ y^2=x^3+ax^2+x.
\]
It is the Kummer cover $w^3=x$, and
$\operatorname{div}(x)=2T-2O$ with $T=(0,0)$
of order two. The curve $E_{a,1}$ is $2$-isogenous
to the original $E$. At the place above two,
pass to a finite extension over which it has good
reduction. The special elliptic fibre is
supersingular. Its geometric two-torsion is
trivial, so $T$ and $O$ meet in the special fibre.
The stable two-pointed model is an elliptic
component joined to a rational component bearing
the two marks. A regular model may contain a
chain of intervening unmarked rational components;
contracting them replaces the chain by a node
of strictly positive thickness.

Normalize this marked model in $w^3=x$.
After a finite extension, all vertical Kummer
valuations are divisible by three. The remaining
normalization is tame, and its local equations
are cube roots of monomials in the marked and
nodal parameters. On the elliptic component,
the node is the only possible branch point.
The degree-zero divisor condition modulo three
excludes branching at just one point. On the
rational component, the two marked orders are
$2,-2$ modulo three, and its node is unramified.
Its connected cyclic cubic cover has genus zero
and three points above the node. At an unramified
node the strict local normalization has three
ordinary nodal branches. Any regular subdivisions
add only rational vertices.

There are two possible covers above the elliptic
component: one connected unramified cubic isogeny,
or three disjoint copies. The connected case has
three edges joining two components, hence graph
first Betti number two. It is excluded by the
Jacobian. A.4 gives
$J_{C_a}\sim E\times F_a^2$, where
$F_a:u^2=t^3-3t+a$ has $j=-216$.
Both factors have potentially good reduction at
two, so the Jacobian has good reduction after a
finite extension. Its torus rank, equal to the
first Betti number of the stable dual graph,
is zero. The split case is therefore forced:
three elliptic tails attached to one rational
center bearing both cusps. The rotation permutes
the tails and preserves their positive arm
lengths. The degree-one maps of the tails to
the good elliptic quotient identify their
symplectic homology spaces.

\paragraph{The actual separating derivation.}
The hyperelliptic involution acts by $+1$ on the
boundary class $t$ and by $-1$ on proper homology.
It therefore gives the equivariant rational
splitting $H_1(X_a)=\mathbf Q_\ell t\oplus H_1(C_a)$
used here. On the three tails it is the elliptic
involution, so their homology classes lie in this
same negative eigenspace. This identifies the
component generators with actual open-curve
homology without adding unspecified boundary
components.
Choose a finite number field $L$ defining the
cover, its deck action and the required data,
with $\Gamma=\overline{[G_L,G_L]}\subset\Delta_0$,
as in A.32. For odd $\ell$, a local tame inertia
element acts trivially on homology after the
preceding semistable extension. Its commutator
with a local Frobenius belongs to $\Gamma$.
The affine edge-twist theorem, applied at an
ordinary point on the rational center, gives
zero first Johnson degree and second-degree
derivation
\[
 D_0(t)=0,\qquad
 D_0|_{V_i}=m_i\operatorname{ad}(\omega_i),\qquad
 \omega_i=[a_i,b_i],\qquad m_i\ne0.
\]
The $m_i$ are the positive arm lengths times a
common nonzero tame commutator scalar.
This uses the same based affine monodromy
statement as A.27 and A.28.

The actual packet representative has the same
outer action. Its difference is inner. Since
both representatives fix the selected cusp,
the possible degree-one inner term is
$\lambda\operatorname{ad}(t)$. If $\lambda\ne0$,
its kernel on the free rational Lie algebra is
$\mathbf Q_\ell t$. Integral primitivity of
the cusp class allows subtraction of the
corresponding actual cusp power. Any nonidentity
remainder would have a higher leading fixed
class, which this kernel excludes. The separated
lower-central filtration then gives precisely
the inertia fixed subgroup.

If $\lambda=0$, both first-degree actions vanish.
Their degree-two difference is
$\operatorname{ad}(z_2)$ and fixes $t$, forcing
$z_2=0$ by the free-Lie centralizer of $t$.
Thus $D_0$ is the actual second Johnson
derivation. The all-degree first-core argument
of A.27, with three elliptic blocks and the
extra letter $t$, gives
\[
 \ker D_0=\mathbb L_{\mathbf Q_\ell}
                      (t,\omega_1,\omega_2,\omega_3).
\]
The proof applies because each $m_i$ is nonzero;
no equality of their values is required.

\paragraph{The non-CM invariant intersection.}
Put $A=\sum a_i$ and $B=\sum b_i$.
The original elliptic subrepresentation is
$V_E=\langle A,B\rangle$, and its symplectic
orthogonal complement is
\[
 V_F^{\oplus2}
   =\left\langle \delta a_i,\delta b_i:1\le i\le3\right\rangle,
 \qquad
 \delta a_i=a_i-A/3,\quad \delta b_i=b_i-B/3.
\]
The deck group acts by permutations on the
three tail indices. Indeed its action on the
elliptic quotient is either the identity or
translation by $T$, both trivial on homology.
The sum-zero part is its standard rational
representation with multiplicity space
$H_1(F_a)$. Thus the arithmetic action on
that part is the same two-dimensional action
on every pair $\delta a_i,\delta b_i$.

The $\Gamma$-action is trivial on $V_E$ and
Zariski dense in $\mathrm{SL}(H_1(F_a))$ on
the multiplicity space, by the non-CM assertion
and Serre theorem checked in A.32. Every
leading class of a $\Gamma$-fixed group
element is therefore fixed by this algebraic
special linear group. In particular it is
fixed by its symplectic unipotent element
\[
 g(a_i)=a_i,\qquad
 g(b_i)=b_i+\delta a_i.
\]
This element fixes $A,B,t$ and commutes with
the deck permutations. It satisfies
\[
 g(\omega_i)=\omega_i+\frac13[A,\delta a_i].
\]
Let $U=\langle\omega_1,\omega_2,\omega_3\rangle$.
If $g(\sum d_i\omega_i)\in U$, its component
in $\bigwedge^2\langle a_1,a_2,a_3\rangle$
is zero. Hence $[A,\sum d_i\delta a_i]=0$.
The vector in the bracket belongs to the
sum-zero plane, whereas the degree-one
centralizer of $A$ is its line. It is
therefore zero, which means $d_1=d_2=d_3$.
We have proved
\[
                         U\cap gU=\mathbf Q_\ell\Omega.
\]

Finally, decompose associative words by the
occurrences of $t$, and pair the letters
in every intervening even-length block.
On each paired tensor block the intersection
of the two subalgebras is the tensor algebra
on $U\cap gU$. This is the same exact
tensor-intersection argument as in A.27.
It gives
\[
 \mathbb L(t,\omega_1,\omega_2,\omega_3)
       \cap g\mathbb L(t,\omega_1,\omega_2,\omega_3)
                      =\mathbb L(t,\Omega).
\]
An actual leading fixed class belongs to the
first algebra by monodromy, and to its
conjugate because it is fixed by $g$.
This proves (A.33.2), including $\ell=3$:
all divisions by three here occur in the
rational graded argument, not in an asserted
integral group operation. The full pro-$2$
result was obtained independently by the
actual induced-topology embedding.

The exact program
\texttt{code/a\_nonstable\_C3\_fixed.py} verifies
the curve and quotient identities, the
symplectic matrix, its deck commutation,
and the rank-five sum of the two
three-dimensional boundary spaces.
Its JSON and stdout are retained.
\end{proof}
